% !TEX root = ../LosningsforslagTRES0312.tex
\chapter{Elektrisitet}
\label{LF:Chapter:elektrisitet}
Dette kapittelet inneholder løsningsforslag til oppgavene i Kapittel~\ref{komp-Chapter:elektrisitet}, \textit{Elektrisitet}, i kompendiet.

\oppgaver{Lading av mobil}{k9oppg1}{%
%
En mobil lades med strømmen $\SI{1.5}{\ampere}$ i $\SI{2.0}{\hour}$. Laderen har spenningen $\SI{5.0}{\volt}$.
\begin{enumerate}[a)]
    \item Hvor mye ladning har gått gjennom ladekabelen?
    \item Hvor mange elektroner tilsvarer det?
    \item Hvor mye energi har laderen levert?
\end{enumerate}
%
%
}

\svar[title={Løsningsforslag}]{%
\noindent
\textbf{a)} Tiden er $t = 2.0\cdot\SI{3600}{\second} = \SI{7200}{\second}$:
\[
  q = I\cdot t = \SI{1.5}{\ampere}\cdot\SI{7200}{\second} = \SI{1.08e4}{\coulomb} \approx \doubleunderline{\SI{1.1e4}{\coulomb}}
\]

\smallskip\noindent
\textbf{b)}
\[
  N = \frac{q}{e} = \frac{\SI{1.08e4}{\coulomb}}{\SI{1.602e-19}{\coulomb}} = \doubleunderline{\num{6.7e22}}
\]

\smallskip\noindent
\textbf{c)}
\[
  W = U\cdot q = \SI{5.0}{\volt}\cdot\SI{1.08e4}{\coulomb} = \doubleunderline{\SI{5.4e4}{\joule}}
\]
}

\oppgaver{Måle strøm og spenning}{k9oppg2}{%
%
Et batteri er koblet til to lyspærer i serie.
\begin{enumerate}[a)]
    \item Tegn kretsen med et amperemeter som måler strømmen fra batteriet, og et voltmeter som måler spenningen over den ene lyspæra.
    \item Hvorfor må amperemeteret kobles i serie, mens voltmeteret kobles i parallell?
    \item Et godt amperemeter har svært liten resistans, mens et godt voltmeter har svært stor resistans. Prøv å forklare hvorfor.
\end{enumerate}
%
%
}

\svar[title={Løsningsforslag}]{%
\noindent
\textbf{a)}
\begin{center}
\begin{circuitikz}
  \draw (0,0) to[battery1, l=$U$] (0,3)
        to[ammeter, l=$I$] (3,3)
        to[lamp] (5.5,3) -- (5.5,1.5) to[lamp] (5.5,0) -- (0,0);
  \draw (5.5,1.5) ++(0,0) -- (7,1.5);
  \draw (7,1.5) to[voltmeter] (7,0) -- (5.5,0);
  \fill (5.5,1.5) circle (2pt) (5.5,0) circle (2pt);
\end{circuitikz}
\end{center}
Amperemeteret står i serie med resten av kretsen. Voltmeteret er koblet over den nederste lyspæra.

\smallskip\noindent
\textbf{b)} Amperemeteret skal måle strømmen \emph{gjennom} kretsen, så strømmen må gå gjennom det. Da må det kobles i serie. Voltmeteret skal måle spenningen \emph{mellom} de to endene av pæra, så det må kobles til begge endene, altså i parallell.

\smallskip\noindent
\textbf{c)} Amperemeteret står i serie. Har det stor resistans, blir strømmen i kretsen mindre, og vi måler noe annet enn det som var der før. Derfor må resistansen være svært liten. Voltmeteret står i parallell. Har det liten resistans, går det mye strøm gjennom voltmeteret i stedet for gjennom pæra, og kretsen endres. Derfor må resistansen være svært stor, slik at nesten ingen strøm går gjennom det.
}

\oppgaver{Ohms lov i tre varianter}{k9oppg3}{%
%
Bruk Ohms lov.
\begin{enumerate}[a)]
    \item En strøm på $\SI{0.050}{\ampere}$ går gjennom en motstand på $\SI{220}{\ohm}$. Hva er spenningen over motstanden?
    \item En motstand på $\SI{4.0}{\ohm}$ kobles til et batteri på $\SI{12}{\volt}$. Hvor stor blir strømmen?
    \item Et batteri på $\SI{1.5}{\volt}$ gir strømmen $\SI{0.30}{\ampere}$ gjennom en liten lyspære. Hva er resistansen i lyspæra?
\end{enumerate}
%
%
}

\svar[title={Løsningsforslag}]{%
\noindent
\textbf{a)} $U = RI = \SI{220}{\ohm}\cdot\SI{0.050}{\ampere} = \doubleunderline{\SI{11}{\volt}}$

\smallskip\noindent
\textbf{b)} $I = \dfrac{U}{R} = \dfrac{\SI{12}{\volt}}{\SI{4.0}{\ohm}} = \doubleunderline{\SI{3.0}{\ampere}}$

\smallskip\noindent
\textbf{c)} $R = \dfrac{U}{I} = \dfrac{\SI{1.5}{\volt}}{\SI{0.30}{\ampere}} = \doubleunderline{\SI{5.0}{\ohm}}$
}

\oppgaver{Måleserie for en motstand}{k9oppg4}{%
%
En elev måler strøm og spenning for en motstand og får disse verdiene:
\begin{center}
\begin{tabular}{c|cccc}
  $I$ / \si{\ampere} & 0.10 & 0.20 & 0.30 & 0.40 \\
  \hline
  $U$ / \si{\volt} & 2.2 & 4.4 & 6.5 & 8.8
\end{tabular}
\end{center}
\begin{enumerate}[a)]
    \item Følger motstanden Ohms lov? Begrunn svaret.
    \item Finn resistansen.
    \item Hvilken spenning trengs for å få strømmen $\SI{0.25}{\ampere}$?
\end{enumerate}
%
%
}

\svar[title={Løsningsforslag}]{%
\noindent
\textbf{a)} Vi regner ut $U/I$ for hver måling: $\SI{22}{\ohm}$, $\SI{22}{\ohm}$, $\SI{21.7}{\ohm}$ og $\SI{22}{\ohm}$. Forholdet er konstant (innenfor måleusikkerheten), så spenningen er proporsjonal med strømmen. Motstanden \doubleunderline{følger Ohms lov}.

\smallskip\noindent
\textbf{b)} Resistansen er $R = U/I \approx \doubleunderline{\SI{22}{\ohm}}$. Det er også stigningstallet til grafen hvis vi plotter $U$ mot $I$.

\smallskip\noindent
\textbf{c)} $U = RI = \SI{22}{\ohm}\cdot\SI{0.25}{\ampere} = \doubleunderline{\SI{5.5}{\volt}}$
}

\oppgaver{Julelys i serie}{k9oppg5}{%
%
Et gammeldags julelys består av 20 like lyspærer koblet i serie. Lenka kobles til $\SI{230}{\volt}$, og strømmen er $\SI{0.10}{\ampere}$.
\begin{enumerate}[a)]
    \item Hva er spenningen over hver pære?
    \item Hva er resistansen i hver pære?
    \item Hvorfor slukner hele lenka når én pære går i stykker?
\end{enumerate}
%
%
}

\svar[title={Løsningsforslag}]{%
\noindent
\textbf{a)} I en seriekobling deles spenningen. Pærene er like, så hver får like mye:
\[
  U_1 = \frac{\SI{230}{\volt}}{20} = \doubleunderline{\SI{11.5}{\volt}}
\]

\smallskip\noindent
\textbf{b)} Den samme strømmen går gjennom alle pærene:
\[
  R_1 = \frac{U_1}{I} = \frac{\SI{11.5}{\volt}}{\SI{0.10}{\ampere}} = \SI{115}{\ohm} \approx \doubleunderline{\SI{1.2e2}{\ohm}}
\]

\smallskip\noindent
\textbf{c)} I en seriekobling har strømmen bare én vei. Når én pære går, blir kretsen \emph{åpen}, og det går ingen strøm noe sted. Da slukner alle pærene.
}

\oppgaver{Lamper i parallell}{k9oppg6}{%
%
To like lamper med resistans $\SI{530}{\ohm}$ hver er koblet i parallell til en stikkontakt med spenningen $\SI{230}{\volt}$.
\begin{enumerate}[a)]
    \item Hvor stor er strømmen gjennom hver lampe?
    \item Hvor stor er den totale strømmen fra stikkontakten?
    \item Hva er den totale resistansen?
    \item Vi kobler til en tredje lampe i parallell. Hva skjer med strømmen gjennom de to første lampene, og med den totale strømmen?
\end{enumerate}
%
%
}

\svar[title={Løsningsforslag}]{%
\noindent
\textbf{a)} Begge lampene har hele spenningen over seg:
\[
  I_1 = \frac{U}{R} = \frac{\SI{230}{\volt}}{\SI{530}{\ohm}} = \SI{0.434}{\ampere} \approx \doubleunderline{\SI{0.43}{\ampere}}
\]

\smallskip\noindent
\textbf{b)} Strømmen deler seg, så $I = I_1 + I_2 = 2\cdot\SI{0.434}{\ampere} = \SI{0.868}{\ampere} \approx \doubleunderline{\SI{0.87}{\ampere}}$.

\smallskip\noindent
\textbf{c)}
\[
  R_\text{tot} = \frac{R_1R_2}{R_1 + R_2} = \frac{\SI{530}{\ohm}\cdot\SI{530}{\ohm}}{\SI{1060}{\ohm}} = \doubleunderline{\SI{265}{\ohm}}
\]
Kontroll: $U/R_\text{tot} = \SI{230}{\volt}/\SI{265}{\ohm} = \SI{0.868}{\ampere}$.

\smallskip\noindent
\textbf{d)} Spenningen over de to første lampene er fortsatt $\SI{230}{\volt}$, så strømmen gjennom dem er \doubleunderline{uendret}. Den nye lampa trekker også $\SI{0.43}{\ampere}$, så den totale strømmen \doubleunderline{øker} til omtrent $\SI{1.3}{\ampere}$. Den totale resistansen blir mindre.
}

\oppgaver{Krets med fire motstander}{k9oppg7}{%
%
Et batteri har polspenningen $\SI{24}{\volt}$. Motstandene $R_1 = \SI{2.0}{\ohm}$ og $R_2 = \SI{4.0}{\ohm}$ er koblet i serie med en parallellkobling av $R_3 = \SI{10}{\ohm}$ og $R_4 = \SI{15}{\ohm}$.
\begin{center}
\begin{circuitikz}
  \draw (0,0) to[battery1, l=$\SI{24}{\volt}$] (0,3)
        to[R, l=$R_1$] (2.5,3) to[R, l=$R_2$] (5,3) -- (7.5,3);
  \draw (5,3) to[R, l_=$R_3$] (5,0);
  \draw (7.5,3) to[R, l=$R_4$] (7.5,0);
  \draw (7.5,0) -- (0,0);
  \fill (5,3) circle (2pt) (5,0) circle (2pt);
\end{circuitikz}
\end{center}
\begin{enumerate}[a)]
    \item Vis at den totale resistansen er $\SI{12}{\ohm}$.
    \item Finn strømmen fra batteriet.
    \item Finn spenningen over parallellkoblingen.
    \item Finn strømmene $I_3$ og $I_4$ gjennom $R_3$ og $R_4$. Kontroller svaret med Kirchhoffs strømlov.
\end{enumerate}
%
%
}

\svar[title={Løsningsforslag}]{%
\noindent
\textbf{a)} Parallellkoblingen:
\[
  R_p = \frac{R_3R_4}{R_3 + R_4} = \frac{\SI{10}{\ohm}\cdot\SI{15}{\ohm}}{\SI{25}{\ohm}} = \SI{6.0}{\ohm}
\]
Den er i serie med $R_1$ og $R_2$: $R_\text{tot} = \SI{2.0}{\ohm} + \SI{4.0}{\ohm} + \SI{6.0}{\ohm} = \doubleunderline{\SI{12}{\ohm}}$.

\smallskip\noindent
\textbf{b)} $I = \dfrac{U}{R_\text{tot}} = \dfrac{\SI{24}{\volt}}{\SI{12}{\ohm}} = \doubleunderline{\SI{2.0}{\ampere}}$

\smallskip\noindent
\textbf{c)} $U_p = R_p I = \SI{6.0}{\ohm}\cdot\SI{2.0}{\ampere} = \doubleunderline{\SI{12}{\volt}}$

\smallskip\noindent
\textbf{d)}
\[
  I_3 = \frac{U_p}{R_3} = \frac{\SI{12}{\volt}}{\SI{10}{\ohm}} = \doubleunderline{\SI{1.2}{\ampere}}, \qquad
  I_4 = \frac{U_p}{R_4} = \frac{\SI{12}{\volt}}{\SI{15}{\ohm}} = \doubleunderline{\SI{0.80}{\ampere}}
\]
Strømloven: $I_3 + I_4 = \SI{1.2}{\ampere} + \SI{0.80}{\ampere} = \SI{2.0}{\ampere} = I$.
}

\oppgaver{Spenningsloven}{k9oppg8}{%
%
To motstander, $R_1$ og $R_2$, er koblet i serie til et batteri på $\SI{9.0}{\volt}$. Et voltmeter over $R_1$ viser $\SI{3.5}{\volt}$.
\begin{enumerate}[a)]
    \item Hva er spenningen over $R_2$? Hvilken av Kirchhoffs lover bruker du?
    \item Resistansen $R_1$ er $\SI{7.0}{\ohm}$. Hvor stor er strømmen i kretsen?
    \item Finn $R_2$.
\end{enumerate}
%
%
}

\svar[title={Løsningsforslag}]{%
\noindent
\textbf{a)} Etter Kirchhoffs \emph{spenningslov} er summen av spenningsfallene lik batterispenningen:
\[
  U_2 = \SI{9.0}{\volt} - \SI{3.5}{\volt} = \doubleunderline{\SI{5.5}{\volt}}
\]

\smallskip\noindent
\textbf{b)} $I = \dfrac{U_1}{R_1} = \dfrac{\SI{3.5}{\volt}}{\SI{7.0}{\ohm}} = \doubleunderline{\SI{0.50}{\ampere}}$

\smallskip\noindent
\textbf{c)} Den samme strømmen går gjennom $R_2$:
\[
  R_2 = \frac{U_2}{I} = \frac{\SI{5.5}{\volt}}{\SI{0.50}{\ampere}} = \doubleunderline{\SI{11}{\ohm}}
\]
}

\oppgaver{Lyspære}{k9oppg9}{%
%
En gammeldags lyspære er merket $\SI{60}{\watt}$ og $\SI{230}{\volt}$.
\begin{enumerate}[a)]
    \item Hvor stor er strømmen gjennom pæra?
    \item Hva er resistansen i pæra når den lyser?
    \item Pæra lyser $\SI{5.0}{\hour}$ hvert døgn i 30 døgn. Hvor mye energi bruker den, i kilowattimer?
    \item En LED-pære på $\SI{8.0}{\watt}$ gir like mye lys. Hvor mye energi sparer vi i løpet av de 30 døgnene ved å bytte?
\end{enumerate}
%
%
}

\svar[title={Løsningsforslag}]{%
\noindent
\textbf{a)} $I = \dfrac{P}{U} = \dfrac{\SI{60}{\watt}}{\SI{230}{\volt}} = \SI{0.261}{\ampere} \approx \doubleunderline{\SI{0.26}{\ampere}}$

\smallskip\noindent
\textbf{b)} $R = \dfrac{U^2}{P} = \dfrac{(\SI{230}{\volt})^2}{\SI{60}{\watt}} = \SI{882}{\ohm} \approx \doubleunderline{\SI{8.8e2}{\ohm}}$

\smallskip\noindent
\textbf{c)} Tiden er $t = 30\cdot\SI{5.0}{\hour} = \SI{150}{\hour}$:
\[
  E = P\cdot t = \SI{0.060}{\kilo\watt}\cdot\SI{150}{\hour} = \doubleunderline{\SI{9.0}{kWh}}
\]

\smallskip\noindent
\textbf{d)} LED-pæra bruker $E = \SI{0.0080}{\kilo\watt}\cdot\SI{150}{\hour} = \SI{1.2}{kWh}$. Vi sparer
\[
  \SI{9.0}{kWh} - \SI{1.2}{kWh} = \doubleunderline{\SI{7.8}{kWh}}
\]
}

\oppgaver{Lading av elbil}{k9oppg10}{%
%
En lader for elbil leverer effekten $\SI{22}{\kilo\watt}$ ved spenningen $\SI{400}{\volt}$.
\begin{enumerate}[a)]
    \item Hvor stor er strømmen?
    \item Batteriet i bilen har en kapasitet på $\SI{60}{kWh}$. Hvor lang tid tar det å lade fra \SI{20}{\percent} til \SI{80}{\percent}? Se bort fra tap.
    \item Hvor mye energi er det, målt i joule?
\end{enumerate}
%
%
}

\svar[title={Løsningsforslag}]{%
\noindent
\textbf{a)} $I = \dfrac{P}{U} = \dfrac{\SI{22e3}{\watt}}{\SI{400}{\volt}} = \doubleunderline{\SI{55}{\ampere}}$

\smallskip\noindent
\textbf{b)} Fra \SI{20}{\percent} til \SI{80}{\percent} er \SI{60}{\percent} av kapasiteten: $E = 0.60\cdot\SI{60}{kWh} = \SI{36}{kWh}$.
\[
  t = \frac{E}{P} = \frac{\SI{36}{kWh}}{\SI{22}{\kilo\watt}} = \SI{1.64}{\hour} \approx \doubleunderline{\SI{1.6}{\hour}}
\]
Det er omtrent 1 time og 40 minutter.

\smallskip\noindent
\textbf{c)} $E = 36\cdot\SI{3.6e6}{\joule} = \SI{1.30e8}{\joule} \approx \doubleunderline{\SI{1.3e8}{\joule}}$
}

\oppgaver{Tre like motstander}{k9oppg11}{%
%
Du har tre like motstander på $\SI{30}{\ohm}$ hver.
\begin{enumerate}[a)]
    \item Hva er den totale resistansen når alle tre er koblet i serie?
    \item Hva er den totale resistansen når alle tre er koblet i parallell?
    \item To av dem kobles i parallell, og den tredje kobles i serie med dem. Hva er den totale resistansen?
    \item Hvilken av koblingene gir størst strøm når de kobles til et batteri på $\SI{12}{\volt}$? Hvor stor er strømmen da?
\end{enumerate}
%
%
}

\svar[title={Løsningsforslag}]{%
\noindent
\textbf{a)} Serie: $R_\text{tot} = 3\cdot\SI{30}{\ohm} = \doubleunderline{\SI{90}{\ohm}}$

\smallskip\noindent
\textbf{b)} Parallell: $\dfrac{1}{R_\text{tot}} = \dfrac{3}{\SI{30}{\ohm}}$, så $R_\text{tot} = \doubleunderline{\SI{10}{\ohm}}$.

\smallskip\noindent
\textbf{c)} To i parallell gir $\SI{15}{\ohm}$. Med den tredje i serie: $R_\text{tot} = \SI{15}{\ohm} + \SI{30}{\ohm} = \doubleunderline{\SI{45}{\ohm}}$.

\smallskip\noindent
\textbf{d)} Størst strøm får vi med minst resistans, altså når alle tre er i \doubleunderline{parallell}:
\[
  I = \frac{\SI{12}{\volt}}{\SI{10}{\ohm}} = \doubleunderline{\SI{1.2}{\ampere}}
\]
}

\oppgaver{Kortslutning}{k9oppg12}{%
%
En mekaniker kommer i skade for å koble polene på et bilbatteri på $\SI{12}{\volt}$ sammen med en skiftenøkkel. Skiftenøkkelen har resistansen $\SI{0.020}{\ohm}$. Se bort fra resistansen i batteriet.
\begin{enumerate}[a)]
    \item Hvor stor blir strømmen?
    \item Hvor stor effekt blir omsatt i skiftenøkkelen?
    \item Hvorfor er kortslutning farlig?
\end{enumerate}
%
%
}

\svar[title={Løsningsforslag}]{%
\noindent
\textbf{a)} $I = \dfrac{U}{R} = \dfrac{\SI{12}{\volt}}{\SI{0.020}{\ohm}} = \doubleunderline{\SI{6.0e2}{\ampere}}$

\smallskip\noindent
\textbf{b)} $P = U\cdot I = \SI{12}{\volt}\cdot\SI{600}{\ampere} = \doubleunderline{\SI{7.2e3}{\watt}}$

\smallskip\noindent
\textbf{c)} Den store strømmen gir svært stor effekt i en liten gjenstand. Skiftenøkkelen og batteripolene blir glovarme på kort tid, og det kan bli gnister, brannskader og brann. Batteriet kan også bli skadet. I virkeligheten har batteriet selv en liten resistans som begrenser strømmen noe, men strømmen blir likevel flere hundre ampere.
}

\oppgaver{Strømregning}{k9oppg13}{%
%
En panelovn har effekten $\SI{1.5}{\kilo\watt}$ og er koblet til $\SI{230}{\volt}$. Den står på $\SI{8.0}{\hour}$ hvert døgn i 30 døgn.
\begin{enumerate}[a)]
    \item Hvor stor er strømmen gjennom ovnen?
    \item Hvor mye energi bruker ovnen i løpet av de 30 døgnene?
    \item Strømprisen er $\SI{1.20}{kr}$ per kilowattime. Hva koster det å ha ovnen på?
\end{enumerate}
%
%
}

\svar[title={Løsningsforslag}]{%
\noindent
\textbf{a)} $I = \dfrac{P}{U} = \dfrac{\SI{1500}{\watt}}{\SI{230}{\volt}} = \SI{6.52}{\ampere} \approx \doubleunderline{\SI{6.5}{\ampere}}$

\smallskip\noindent
\textbf{b)} Tiden er $t = 30\cdot\SI{8.0}{\hour} = \SI{240}{\hour}$:
\[
  E = P\cdot t = \SI{1.5}{\kilo\watt}\cdot\SI{240}{\hour} = \doubleunderline{\SI{3.6e2}{kWh}}
\]

\smallskip\noindent
\textbf{c)} Kostnaden er $\SI{360}{kWh}\cdot\SI{1.20}{kr/kWh} = \SI{432}{kr} \approx \doubleunderline{\SI{4.3e2}{kr}}$.
}

\oppgaver{Krets med voltmeter}{k9oppg14}{%
%
Et batteri på $\SI{12}{\volt}$ er koblet i serie med en motstand $R_1 = \SI{4.0}{\ohm}$ og en lyspære. Et voltmeter koblet over lyspæra viser $\SI{8.0}{\volt}$.
\begin{enumerate}[a)]
    \item Tegn kretsen med voltmeteret.
    \item Hva er spenningen over $R_1$?
    \item Hvor stor er strømmen i kretsen?
    \item Hva er resistansen i lyspæra?
    \item Hvor stor effekt lyser pæra med? Hvor stor effekt leverer batteriet?
\end{enumerate}
%
%
}

\svar[title={Løsningsforslag}]{%
\noindent
\textbf{a)}
\begin{center}
\begin{circuitikz}
  \draw (0,0) to[battery1, l=$\SI{12}{\volt}$] (0,3)
        to[R, l=$R_1$] (3,3)
        to[lamp] (3,0) -- (0,0);
  \draw (3,3) -- (4.6,3) to[voltmeter] (4.6,0) -- (3,0);
\end{circuitikz}
\end{center}

\noindent
\textbf{b)} Spenningsloven: $U_1 = \SI{12}{\volt} - \SI{8.0}{\volt} = \doubleunderline{\SI{4.0}{\volt}}$

\smallskip\noindent
\textbf{c)} $I = \dfrac{U_1}{R_1} = \dfrac{\SI{4.0}{\volt}}{\SI{4.0}{\ohm}} = \doubleunderline{\SI{1.0}{\ampere}}$

\smallskip\noindent
\textbf{d)} $R_\text{L} = \dfrac{\SI{8.0}{\volt}}{\SI{1.0}{\ampere}} = \doubleunderline{\SI{8.0}{\ohm}}$

\smallskip\noindent
\textbf{e)} Pæra: $P_\text{L} = U_\text{L}I = \SI{8.0}{\volt}\cdot\SI{1.0}{\ampere} = \doubleunderline{\SI{8.0}{\watt}}$. Batteriet: $P = \SI{12}{\volt}\cdot\SI{1.0}{\ampere} = \doubleunderline{\SI{12}{\watt}}$. Resten, $\SI{4.0}{\watt}$, blir til varme i $R_1$.
}
